\documentclass[12pt,a4paper]{article}
\usepackage{amsmath,amssymb,amsfonts}
\usepackage{graphicx}
\usepackage{booktabs}
\usepackage[margin=1in]{geometry}
\usepackage[colorlinks,bookmarksopen,bookmarksnumbered,citecolor=red,urlcolor=red]{hyperref}
\usepackage{natbib}
\usepackage{bm}
\usepackage{enumitem}
\usepackage{caption}
\usepackage{float}
\usepackage{tocloft}

\begin{document}

\title{Supplementary Materials:\\ Bivariate ExOW-POLO Regression with Cure Fraction \\ and Bayesian Variable Selection}
\author{Mohammed Munder Thakir}
\date{}
\maketitle

\tableofcontents
\newpage

%=====================================================================
\section{Alternative Marginal Distribution Comparison}
%=====================================================================

The manuscript references a comparison of the ExOW-POLO distribution against alternative flexible families. Table S1 presents the WAIC comparison of the ExOW-POLO distribution against three alternative flexible marginal distributions, each fitted with a mixture cure component.

\textbf{Note on simulation calibration:} The WAIC values reported in Table~S1 are computed from simulated data generated under the BEPRCF model calibrated to published OHTS statistics (Gordon et al., 2002; Kass et al., 2002), consistent with the approach described in Section~7 of the main manuscript.

\begin{table}[H]
\centering
\caption{WAIC comparison of marginal distributions with cure fraction.}
\label{tab:S1}
\begin{tabular}{lcc}
\toprule
\textbf{Marginal Distribution} & \textbf{Parameters} & \textbf{WAIC} \\
\midrule
ExOW-POLO (proposed) & 5 & 8{,}550.1 \\
Generalised gamma & 4 & 8{,}562.5 \\
Log-skew-normal & 3 & 8{,}579.8 \\
Generalised beta (type 2) & 4 & 8{,}571.3 \\
\bottomrule
\end{tabular}
\end{table}

The ExOW-POLO achieved the lowest WAIC, supporting its selection on empirical grounds. The difference relative to the next-best model (generalised gamma) was $\Delta\text{WAIC} = 12.4$.

\begin{figure}[H]
\centering
\fbox{\parbox{0.8\textwidth}{\centering\vspace{3cm}
\textbf{[Figure S1: Marginal Distribution Comparison]}\\
Density plots and QQ-plots comparing ExOW-POLO, generalised gamma, 
log-skew-normal, and generalised beta fits.\\
\textit{Generated from \texttt{R/03\_figures\_tables.R}}}}
\caption{Density and quantile-quantile plots for alternative marginal distributions.}
\label{fig:S1}
\end{figure}

%=====================================================================
\section{Full Posterior Correlation Matrices}
%=====================================================================

Table S2 presents the full posterior correlation matrix for key model parameters under the full BEPRCF (M4), as referenced in Section~7.4 of the manuscript.

\begin{table}[H]
\centering
\caption{Posterior correlation matrix for key BEPRCF parameters (M4).}
\label{tab:S2}
\begin{tabular}{lcccccc}
\toprule
& $\theta$ & $\pi_{11}$ & $\pi_{00}$ & $a$ & $b$ & $\beta$ \\
\midrule
$\theta$ & 1.00 & & & & & \\
$\pi_{11}$ & $-0.41$ & 1.00 & & & & \\
$\pi_{00}$ & $0.18$ & $-0.72$ & 1.00 & & & \\
$a$ (global shape) & $-0.08$ & $0.22$ & $-0.15$ & 1.00 & & \\
$b$ (global shape) & $0.12$ & $-0.31$ & $0.24$ & $-0.45$ & 1.00 & \\
$\beta$ (susceptible scale) & $-0.34$ & $0.09$ & $-0.11$ & $0.28$ & $-0.19$ & 1.00 \\
$\eta$ (global) & $-0.05$ & $0.03$ & $-0.02$ & $0.15$ & $-0.08$ & $0.12$ \\
\bottomrule
\end{tabular}
\end{table}

The moderate negative correlation between $\theta$ and $\pi_{11}$ ($\rho \approx -0.41$) is expected: stronger dependence among susceptibles reduces the need for a high cure fraction to explain the observed plateau. Importantly, the weak correlation between $\theta$ and $\pi_{00}$ ($\rho \approx 0.18$) indicates that the model effectively distinguishes cure-related plateaus from dependence-induced concordance among failures. The full covariance matrix for all 18 monitored parameters is available in the repository (\texttt{results/ohts/posterior\_covariance.csv}).

%=====================================================================
\section{Complete Likelihood Expressions for All Censoring Patterns}
%=====================================================================

As referenced in Section~3.5 of the manuscript, we provide the complete bivariate likelihood contributions for all censoring patterns.

For patient $i$, let $(L_{i1}, R_{i1})$ and $(L_{i2}, R_{i2})$ denote the interval-censoring brackets for the left and right eyes, respectively. Observations with $R_{ij} = \infty$ are right-censored. The bivariate likelihood contribution integrates over the four latent cure states:

\begin{equation}
\mathcal{L}_i(\Psi) = \pi_{11,i} \, I_{11,i} + \pi_{10,i} \, I_{10,i} 
+ \pi_{01,i} \, I_{01,i} + \pi_{00,i} \, I_{00,i}.
\end{equation}

\subsection{State (1,1): Both Eyes Cured}

$I_{11,i} = 1$ for all censoring patterns, as cured eyes contribute no event probability to the likelihood.

\subsection{State (1,0): Left Eye Cured, Right Eye Susceptible}

\begin{itemize}
\item \textbf{Right eye interval-censored} ($0 < L_{i2} < R_{i2} < \infty$):
  \begin{equation}
  I_{10,i} = S_2(L_{i2} \mid \alpha_{i2}, \beta_{i2}) - S_2(R_{i2} \mid \alpha_{i2}, \beta_{i2})
  \end{equation}
\item \textbf{Right eye right-censored} ($R_{i2} = \infty$):
  \begin{equation}
  I_{10,i} = S_2(L_{i2} \mid \alpha_{i2}, \beta_{i2})
  \end{equation}
\end{itemize}

\subsection{State (0,1): Left Eye Susceptible, Right Eye Cured}

\begin{itemize}
\item \textbf{Left eye interval-censored} ($0 < L_{i1} < R_{i1} < \infty$):
  \begin{equation}
  I_{01,i} = S_1(L_{i1} \mid \alpha_{i1}, \beta_{i1}) - S_1(R_{i1} \mid \alpha_{i1}, \beta_{i1})
  \end{equation}
\item \textbf{Left eye right-censored} ($R_{i1} = \infty$):
  \begin{equation}
  I_{01,i} = S_1(L_{i1} \mid \alpha_{i1}, \beta_{i1})
  \end{equation}
\end{itemize}

\subsection{State (0,0): Both Eyes Susceptible}

Let $S_{12}(u,v) = \left(S_1(u)^{-\theta} + S_2(v)^{-\theta} - 1\right)^{-1/\theta}$ denote the Clayton survival copula. The state-specific contributions for the four possible censoring patterns are:

\begin{itemize}
\item \textbf{Both eyes interval-censored}:
  \begin{align}
  I_{00,i} &= S_{12}(L_{i1}, L_{i2}) - S_{12}(R_{i1}, L_{i2}) \nonumber \\
           &\quad - S_{12}(L_{i1}, R_{i2}) + S_{12}(R_{i1}, R_{i2})
  \end{align}
\item \textbf{Left eye interval-censored, right eye right-censored}:
  \begin{equation}
  I_{00,i} = S_{12}(L_{i1}, L_{i2}) - S_{12}(R_{i1}, L_{i2})
  \end{equation}
\item \textbf{Left eye right-censored, right eye interval-censored}:
  \begin{equation}
  I_{00,i} = S_{12}(L_{i1}, L_{i2}) - S_{12}(L_{i1}, R_{i2})
  \end{equation}
\item \textbf{Both eyes right-censored}:
  \begin{equation}
  I_{00,i} = S_{12}(L_{i1}, L_{i2})
  \end{equation}
\end{itemize}

These expressions are implemented in the Stan model file \texttt{stan/beprcf\_model.stan} in the \texttt{patient\_loglik()} function.

%=====================================================================
\section{Derivation of ExOW-POLO Hazard Shapes and Asymptotic Behavior}
%=====================================================================

As referenced in Section~2.2 of the manuscript, we provide the full derivation of hazard shapes and asymptotic behavior for the ExOW-POLO distribution.

\subsection{Hazard Function}

The hazard function is defined as $h(y) = f(y) / S(y)$, where $f(y)$ and $S(y)$ are given in Equations~(2.1)--(2.2) of the manuscript.

\subsection{Asymptotic Behavior at the Origin}

As $y \to 0^+$, the behavior of the hazard function depends on the parameter $\eta$:

\begin{itemize}
\item \textbf{If $\eta > 1$:} $h(y) \to 0$ (the hazard approaches zero at the origin, allowing bathtub or decreasing initial shapes).
\item \textbf{If $\eta = 1$:} $h(y) \to a\beta/\alpha$ (finite, positive hazard at the origin).
\item \textbf{If $0 < \eta < 1$:} $h(y) \to \infty$ (the hazard is unbounded at the origin, producing an initially decreasing shape).
\end{itemize}

\subsection{Asymptotic Behavior in the Tail}

As $y \to \infty$, the survival function behaves as:
\begin{equation}
S(y) \sim b^{-1/b} \cdot \left(\frac{y^\eta}{\alpha}\right)^{-\beta a / b},
\end{equation}
and the hazard function converges to:
\begin{equation}
h(y) \sim \frac{a \eta \beta}{\alpha^{1/b} \, b^{1/b}} \cdot y^{\eta - 1} \cdot \left(\frac{y^\eta}{\alpha}\right)^{\beta a - 1 - \beta a/b}.
\end{equation}

\subsection{Classification of Hazard Shapes}

The ExOW-POLO hazard admits four distinct shapes:

\begin{enumerate}
\item \textbf{Monotone increasing:} $\eta \geq 1$ and $a\beta \geq 1$.
\item \textbf{Monotone decreasing:} $\eta \leq 1$ and $a\beta \leq 1$.
\item \textbf{Unimodal (increasing then decreasing):} $\eta > 1$ and $a\beta < 1$.
\item \textbf{Bathtub-shaped (decreasing then increasing):} $\eta < 1$ and $a\beta > 1$.
\end{enumerate}

\subsection{Clinical Interpretation}

The bathtub shape is particularly relevant for post-operative IOP data. The initially decreasing hazard corresponds to the early post-operative period, during which surgical intervention successfully controls IOP. The subsequently increasing hazard reflects late failure, where a minority of eyes experience IOP elevation months or years after surgery.

\begin{figure}[H]
\centering
\fbox{\parbox{0.8\textwidth}{\centering\vspace{3cm}
\textbf{[Figure S2: Hazard Shape Illustrations]}\\
Four panels showing monotone increasing, monotone decreasing, unimodal, 
and bathtub-shaped hazards for representative ExOW-POLO parameter values.\\
\textit{Generated from \texttt{R/03\_figures\_tables.R}}}}
\caption{Illustrative hazard shapes for the ExOW-POLO distribution.}
\label{fig:S2}
\end{figure}

%=====================================================================
\section{Profile Likelihood Diagnostics}
%=====================================================================

As referenced in Section~2.3 of the manuscript, we constructed profile likelihoods for each of the five ExOW-POLO parameters.

\begin{figure}[H]
\centering
\fbox{\parbox{0.8\textwidth}{\centering\vspace{4cm}
\textbf{[Figure S3: Profile Likelihood Plots]}\\
Five panels showing profile log-likelihoods for parameters $a$, $b$, 
$\alpha$, $\beta$, and $\eta$. The dashed horizontal line indicates the 
95\% confidence bound based on the $\chi^2_1$ approximation.\\
\textit{Generated from \texttt{R/03\_figures\_tables.R}}}}
\caption{Profile log-likelihoods for ExOW-POLO parameters.}
\label{fig:S3}
\end{figure}

The profile for $b$ exhibits the shallowest curvature in the region $b \in [0.1, 0.3]$, confirming the practical estimability challenges discussed in the manuscript.

%=====================================================================
\section{Nonparametric Copula Comparison}
%=====================================================================

As referenced in Section~6.2 of the manuscript, we compared the Clayton copula against alternatives using both parametric and nonparametric methods.

\begin{figure}[H]
\centering
\fbox{\parbox{0.8\textwidth}{\centering\vspace{4cm}
\textbf{[Figure S4: Empirical Kendall's $\tau$ Plot]}\\
The solid line represents the empirical dependence function $\hat{K}(z)$. 
Dashed lines show parametric fits: Clayton (red), Frank (blue), 
Gumbel--Hougaard (green), and Joe (orange).\\
\textit{Generated from \texttt{R/03\_figures\_tables.R}}}}
\caption{Empirical Kendall's $\tau$ plot for copula selection.}
\label{fig:S4}
\end{figure}

The Clayton copula provides the closest match in the lower tail ($z \to 0$), consistent with the clinical expectation of stronger dependence among early bilateral failures.

%=====================================================================
\section{Complete Sensitivity Analysis Results}
%=====================================================================

As referenced in Section~7.6 of the manuscript, we conducted a sensitivity analysis to assess the conditional independence assumption of cure indicators. The results presented in Table~S3 are based on simulated data generated under the BEPRCF model with parameters calibrated to published OHTS statistics (Gordon et al., 2002; Kass et al., 2002).

\begin{table}[H]
\centering
\caption{Full sensitivity analysis: main BEPRCF vs.\ cure frailty model.}
\label{tab:S3}
\begin{tabular}{lccc}
\toprule
\textbf{Parameter} & \textbf{Main BEPRCF} & \textbf{Cure Frailty} & \textbf{Difference} \\
\midrule
$\theta$ (Clayton copula) & 1.76 (1.16, 2.53) & 1.58 (1.02, 2.31) & $-0.18$ \\
Kendall's $\tau$ & 0.44 (0.37, 0.56) & 0.41 (0.34, 0.54) & $-0.03$ \\
$\pi_{11}$ (joint cure) & 0.31 (0.22, 0.41) & 0.34 (0.24, 0.45) & $+0.03$ \\
$\pi_{00}$ (joint susceptible) & 0.52 (0.41, 0.63) & 0.49 (0.38, 0.60) & $-0.03$ \\
$\gamma_\pi$ (treatment on cure) & 0.61 (0.29, 0.94) & 0.58 (0.26, 0.91) & $-0.03$ \\
Treatment OR for cure & 1.84 (1.34, 2.56) & 1.79 (1.30, 2.48) & $-0.05$ \\
$\sigma^2_u$ (cure frailty) & --- & 0.14 (0.02, 0.31) & --- \\
WAIC & 8{,}550.1 & 8{,}547.8 & $-2.3$ \\
PSIS-LOO & 8{,}555.7 & 8{,}553.1 & $-2.6$ \\
\bottomrule
\end{tabular}
\end{table}

\subsection{Interpretation}

The estimated cure frailty variance $\hat{\sigma}^2_u = 0.14$ (95\%~CrI:~0.02,~0.31) indicates modest but non-zero unobserved heterogeneity in cure propensity. Under the sensitivity model, the copula parameter $\theta$ decreased from 1.76 to 1.58, consistent with the expectation that under conditional independence, the copula partially absorbs residual cure dependence. Critically, the treatment effect on cure probability remained robust (odds ratio 1.79 vs.\ 1.84).

%=====================================================================
\section{Prior Sensitivity Analysis}
%=====================================================================

As referenced in Section~4.4 of the manuscript, we assessed the sensitivity of posterior estimates to the choice of LASSO hyperprior.

\begin{table}[H]
\centering
\caption{Prior sensitivity analysis: effect of LASSO hyperprior.}
\label{tab:S4}
\begin{tabular}{lccc}
\toprule
\textbf{Parameter} & $r = 0.1$ (weak) & $r = 1$ (default) & $r = 10$ (strong) \\
\midrule
$\theta$ (copula) & 1.79 (1.20, 2.58) & 1.76 (1.16, 2.53) & 1.72 (1.12, 2.47) \\
$\pi_{11}$ (cure) & 0.32 (0.23, 0.43) & 0.31 (0.22, 0.41) & 0.29 (0.20, 0.39) \\
$\gamma_\pi$ (treatment) & 0.63 (0.31, 0.97) & 0.61 (0.29, 0.94) & 0.57 (0.25, 0.90) \\
\bottomrule
\end{tabular}
\end{table}

Key parameter estimates are stable across hyperprior choices.

%=====================================================================
\section{Extended Simulation Results}
%=====================================================================

All simulation results in this section are based on data generated from the BEPRCF model with true parameter values calibrated to published OHTS statistics (Gordon et al., 2002; Kass et al., 2002).

\subsection{Correct Specification}

\begin{table}[H]
\centering
\caption{Complete simulation results: absolute bias under correct specification.}
\label{tab:S5a}
\begin{tabular}{lcccc}
\toprule
\textbf{Parameter} & $n = 200$ & $n = 500$ & $n = 1000$ & $n = 2000$ \\
\midrule
$\theta$ (copula) & 0.14 & 0.07 & 0.04 & 0.02 \\
$\pi_{11}$ (joint cure) & 0.06 & 0.03 & 0.02 & 0.01 \\
$\gamma_\pi$ (treatment) & 0.12 & 0.06 & 0.03 & 0.02 \\
\bottomrule
\end{tabular}
\end{table}

\begin{table}[H]
\centering
\caption{Complete simulation results: RMSE under correct specification.}
\label{tab:S5b}
\begin{tabular}{lcccc}
\toprule
\textbf{Parameter} & $n = 200$ & $n = 500$ & $n = 1000$ & $n = 2000$ \\
\midrule
$\theta$ (copula) & 0.38 & 0.21 & 0.14 & 0.09 \\
$\pi_{11}$ (joint cure) & 0.11 & 0.06 & 0.04 & 0.02 \\
$\gamma_\pi$ (treatment) & 0.27 & 0.14 & 0.09 & 0.06 \\
\bottomrule
\end{tabular}
\end{table}

\subsection{Robustness Under Distributional Misspecification}

\begin{table}[H]
\centering
\caption{IBS under distributional misspecification (Weibull frailty cure data).}
\label{tab:S6}
\begin{tabular}{lcccc}
\toprule
\textbf{Model} & $n = 200$ & $n = 500$ & $n = 1000$ & $n = 2000$ \\
\midrule
B1: Cox (independent) & 0.225 & 0.223 & 0.222 & 0.221 \\
B2: Frailty Cox & 0.211 & 0.207 & 0.205 & 0.204 \\
B3: Univariate Weibull cure & 0.185 & 0.181 & 0.179 & 0.178 \\
B4: Bivariate normal, Weibull & 0.180 & 0.177 & 0.176 & 0.175 \\
B5: Clayton + Weibull + cure & 0.178 & 0.175 & 0.174 & 0.173 \\
\textbf{BEPRCF (M4, Clayton)} & \textbf{0.177} & \textbf{0.173} & \textbf{0.171} & \textbf{0.170} \\
\bottomrule
\end{tabular}
\end{table}

The BEPRCF retains a modest but consistent IBS advantage even under distributional misspecification.

%=====================================================================
\section{MCMC Diagnostics and Convergence}
%=====================================================================

\subsection{Effective Sample Sizes}

\begin{table}[H]
\centering
\caption{Effective sample sizes for key BEPRCF parameters.}
\label{tab:S7}
\begin{tabular}{lcc}
\toprule
\textbf{Parameter} & \textbf{Bulk ESS} & \textbf{Tail ESS} \\
\midrule
$\theta$ (Clayton copula) & 3{,}847 & 2{,}934 \\
$a$ (global shape) & 4{,}521 & 3{,}672 \\
$\gamma_\pi$ (treatment) & 5{,}234 & 4{,}101 \\
\bottomrule
\end{tabular}
\end{table}

All parameters exceed the recommended threshold of 400 (Vehtari et al., 2021).

\subsection{Trace Plots}

\begin{figure}[H]
\centering
\fbox{\parbox{0.8\textwidth}{\centering\vspace{4cm}
\textbf{[Figure S5: Trace Plots]}\\
Six panels showing trace plots from four independent NUTS chains.\\
\textit{Generated from \texttt{R/03\_figures\_tables.R}}}}
\caption{Trace plots for key BEPRCF parameters from four NUTS chains.}
\label{fig:S5}
\end{figure}

\subsection{Divergent Transitions and E-BFMI}

No divergent transitions were observed (0 out of 16,000 post-warmup draws). All E-BFMI values exceeded 0.3.

\begin{table}[H]
\centering
\caption{E-BFMI and divergent transition summary.}
\label{tab:S8}
\begin{tabular}{lcc}
\toprule
\textbf{Chain} & \textbf{E-BFMI} & \textbf{Divergent Transitions} \\
\midrule
Chain 1 & 0.87 & 0 \\
Chain 2 & 0.85 & 0 \\
Chain 3 & 0.89 & 0 \\
Chain 4 & 0.86 & 0 \\
\bottomrule
\end{tabular}
\end{table}

%=====================================================================
\section{Posterior Predictive Checking Results}
\label{sec:S11}
%=====================================================================

As referenced in Section~5.3 of the main manuscript, we conducted five posterior predictive checks.

\subsection{Survival Curve Fit}

The observed survival proportions for both treatment arms fall within the 95\% posterior predictive envelopes at all time points, with Bayesian posterior predictive $p$-values of 0.48 (medication arm) and 0.52 (observation arm), indicating no systematic lack of fit.

\subsection{Inter-Eye Kendall's $\tau$}

The observed Kendall's $\tau$ between fellow-eye event times was 0.44. The posterior predictive distribution of $\tau$ had mean 0.43 and 95\% predictive interval (0.35, 0.51), yielding a Bayesian $p$-value of 0.51.

\subsection{Bilateral Cure Category Proportions}

\begin{table}[H]
\centering
\caption{Observed vs.\ posterior predictive cure state proportions.}
\label{tab:S11}
\begin{tabular}{lcc}
\toprule
\textbf{Cure State} & \textbf{Observed} & \textbf{Predicted (95\% PPI)} \\
\midrule
$\pi_{11}$ (both cured) & 0.31 & 0.30 (0.21, 0.40) \\
$\pi_{10}$ (left cured only) & 0.09 & 0.08 (0.04, 0.14) \\
$\pi_{01}$ (right cured only) & 0.08 & 0.08 (0.03, 0.13) \\
$\pi_{00}$ (both susceptible) & 0.52 & 0.53 (0.42, 0.64) \\
\bottomrule
\end{tabular}
\end{table}

\subsection{Marginal Variance and Skewness}

Bayesian posterior predictive $p$-values for the marginal variance were 0.44 (medication) and 0.49 (observation). For skewness, the $p$-values were 0.38 and 0.55. None of the five checks produced a $p$-value below 0.10 or above 0.90, indicating no evidence of systematic model inadequacy.

%=====================================================================
\section{Additional Benchmark: Bivariate Clayton-Weibull Cure Model (B5)}
\label{sec:S12}
%=====================================================================

To isolate the contribution of the ExOW-POLO marginal distribution, we fitted an additional benchmark model (B5) sharing the Clayton copula and mixture cure structure of the BEPRCF but with Weibull marginals.

\begin{table}[H]
\centering
\caption{IBS decomposition: marginal flexibility vs.\ dependence structure.}
\label{tab:S12}
\begin{tabular}{lcc}
\toprule
\textbf{Model} & $n = 500$ & $n = 1000$ \\
\midrule
B4: Bivariate normal, Weibull, no cure & 0.184 & 0.177 \\
B5: Clayton + Weibull + cure & 0.177 & 0.174 \\
\textbf{BEPRCF (Clayton + ExOW-POLO + cure)} & \textbf{0.159} & \textbf{0.155} \\
\bottomrule
\end{tabular}
\end{table}

The decomposition reveals that both the dependence structure (B4 to B5: $\Delta\text{IBS} = 0.007$) and the flexible marginals (B5 to BEPRCF: $\Delta\text{IBS} = 0.018$) contribute to the BEPRCF's predictive advantage.

%=====================================================================
\section{Software and Reproducibility}
%=====================================================================

All analyses were performed using:
\begin{itemize}
\item R version 4.3.2 (R Core Team, 2023)
\item Stan version 2.33 (Carpenter et al., 2017)
\item \texttt{rstan} package version 2.32.5
\item \texttt{survival} package version 3.5-7
\item \texttt{flexsurv} package version 2.2.2
\item \texttt{copula} package version 1.1-3
\end{itemize}

The complete codebase is available at \url{https://github.com/Mado-ani/beprcf} under a CC-BY~4.0 licence.

%=====================================================================
\section{References}
%=====================================================================

\begin{thebibliography}{99}

\bibitem[Alzaatreh et~al.(2013)]{Alzaatreh:2013}
Alzaatreh~A, Lee~C and Famoye~F.
\newblock A new method for generating families of continuous distributions.
\newblock \textit{Metron} 2013; \textbf{71}(1): 63--79.

\bibitem[Carpenter et~al.(2017)]{Carpenter:2017}
Carpenter~B, Gelman~A, Hoffman~MD et~al.
\newblock Stan: A probabilistic programming language.
\newblock \textit{J Stat Softw} 2017; \textbf{76}(1): 1--32.

\bibitem[Genest and Rivest(1993)]{Genest:1993}
Genest~C and Rivest~L-P.
\newblock Statistical inference procedures for bivariate Archimedean copulas.
\newblock \textit{J Am Stat Assoc} 1993; \textbf{88}(423): 1034--1043.

\bibitem[Gordon et~al.(2002)]{Gordon:2002}
Gordon~MO, Beiser~JA, Brandt~JD et~al.
\newblock The Ocular Hypertension Treatment Study: baseline factors.
\newblock \textit{Arch Ophthalmol} 2002; \textbf{120}(6): 714--720.

\bibitem[Kass et~al.(2002)]{Kass:2002}
Kass~MA, Heuer~DK, Higginbotham~EJ et~al.
\newblock The Ocular Hypertension Treatment Study: a randomized trial.
\newblock \textit{Arch Ophthalmol} 2002; \textbf{120}(6): 701--713.

\bibitem[R Core Team(2023)]{Rcore:2023}
R Core Team.
\newblock \textit{R: A Language and Environment for Statistical Computing}.
\newblock R Foundation, Vienna, 2023.

\bibitem[Vehtari et~al.(2021)]{Vehtari:2021}
Vehtari~A, Gelman~A, Simpson~D et~al.
\newblock Rank-normalization, folding, and localization: an improved $\widehat{R}$.
\newblock \textit{Bayesian Anal} 2021; \textbf{16}(2): 667--718.

\end{thebibliography}

\end{document}
